\documentclass[11pt]{article}

\usepackage[margin=1in]{geometry}
\usepackage{amsmath,amssymb,amsthm}
\usepackage{enumitem}
\usepackage{hyperref}
\usepackage[T1]{fontenc}

\setlength{\parindent}{0pt}
\setlength{\parskip}{0.75em}

\begin{document}

\begin{center}
    {\LARGE \textbf{Qualifying Exam Syllabus}}\\[1em]
\end{center}

\textbf{Name:} Songyu Ye \\
\textbf{Date:} May 14, 2026 \\
\textbf{Time:} 9am - 12pm \\
\textbf{Location:} Evans 732 \\
\textbf{Committee:} Constantin Teleman (advisor), David Nadler (exam chair), Hannah Larson, Richard Borcherds

\vspace{1em}

\section*{1. Major Topic: Representation Theory (Algebra)}

\begin{itemize}[leftmargin=2em]
    \item \textbf{Lie algebras:} solvable, nilpotent, and semisimple Lie algebras; Cartan subalgebras; root systems; Weyl groups; Dynkin diagrams; classification of simple Lie algebras; Killing form; complete reducibility; examples: $\mathfrak{sl}_n$, $\mathfrak{so}_n$, $\mathfrak{sp}_{2n}$, $\mathfrak{u}_n$.

    \item \textbf{Lie groups and algebraic groups:} fundamental group and center of Lie groups; simply connected and adjoint forms associated to a simple Lie algebra; real forms, compact forms, split forms; Cartan involution; maximal compact subgroups of real Lie groups; reductive, Borel subgroups, parabolic subgroups, Levi subgroups.

    \item \textbf{Representation theory:} universal enveloping algebra; PBW theorem; Verma modules; highest weight modules; Casimir elements; central characters; Weyl character formula.
\end{itemize}

\textbf{Reference:} Kirillov, \emph{Introduction to Lie Groups and Lie Algebras}. Varadarajan, \emph{Lie Groups, Lie Algebras, and
Their Representations}.

\section*{2. Major Topic: Algebraic Geometry (Algebra)}

\begin{itemize}[leftmargin=2em]
    \item \textbf{Sheaves:} presheaves and sheaves; morphisms of sheaves; stalks; exact sequences of sheaves; skyscraper sheaves; locally free sheaves; quasi-coherent sheaves.
\item \textbf{Morphisms of schemes:} affine morphisms; open and closed immersions; finite and integral morphisms; separated, proper, and projective morphisms; flat, smooth, unramified, and \'{e}tale morphisms; fiber products and base change. 
    \item \textbf{Divisors and bundles:} Cartier divisors and Weil divisors; linear equivalence; Picard groups; linear systems; line bundles; ampleness and very ampleness; vector bundles; Chern classes; normal bundles of subvarieties; adjunction formula; sheaf of differentials.

    \item \textbf{Sheaf cohomology and derived functors:} cohomology of sheaves; cohomology of projective spaces; pullback and pushforward; higher direct images; vanishing theorems; Serre duality; flatness; cohomology and base change.

    \item \textbf{Curves:} Riemann--Roch; Riemann--Hurwitz; curves in projective spaces; elliptic curves; hyperelliptic curves.
\end{itemize}

\textbf{Reference:} Hartshorne, \emph{Algebraic Geometry}, Chapters 2--4.

\section*{3. Minor Topic: Algebraic Topology (Geometry)}

\begin{itemize}[leftmargin=2em]
    \item \textbf{Fundamental groups and homotopy:} Fundamental group; Van Kampen's theorem; fundamental groups of graphs and surfaces; covering space theory; universal covers; deck transformations; CW complexes; higher homotopy groups; Whitehead's theorem; Hurewicz theorem. 

    \item \textbf{Homology and cohomology:} cellular and simplicial homology; de Rham cohomology; long exact sequences; Mayer--Vietoris sequence; relative homology and cohomology; excision; universal coefficient theorem; K\"unneth theorem; cup product in cohomology; Poincar\'e duality for closed orientable manifolds; 
\end{itemize}

\textbf{Reference:} Hatcher, \emph{Algebraic Topology}, Chapters 1--4, excluding additional topics.



\end{document}